\section{Conclusion}
\label{sec:conclusion}

As thermal constraints increasingly dictate the performance ceilings of modern heterogeneous architectures, relying on simplified thermal abstractions severely handicaps scheduling decisions. In this paper, we made two contributions to address this gap. First, we instrumented gem5 with a \texttt{ThermalNodeStats} patch that directly exports $T_{die}$, $T_{pkg}$, and $T_{hs}$ to simulation statistics, enabling direct multi-node temperature export. Second, using these gem5-calibrated parameters ($\tau_{hs} = \SI{120}{\second}$), we showed that a high-fidelity V3 3-Node Cauer model correctly defers the first thermal intervention by \textbf{85 simulated seconds} compared to a 2-node baseline, because the model captures heatsink inertia that 2-node approximations discard.

Building on the 3-node model, the Phase~6 package-aware scheduling policy reduces DVFS oscillation events from \textbf{13 to 1} over a 600-second sustained workload, lowering modeled peak package temperature by \textbf{4.5\,°C}. This is achieved by prioritizing task migration when $T_{pkg}$ retains headroom, reserving DVFS as a fallback only when heatsink capacity is genuinely exhausted. The trade-off is a 12\% reduction in projected throughput, an acceptable cost for significantly reduced thermal cycling. All results are modeled via trace-driven co-simulation; future work will validate these findings on physical ARM big.LITTLE hardware and explore applying this framework to predictive, ML-based schedulers for emerging 2.5D/3D chiplet architectures.
